\chapter{Перечень продовольственных категорий, составляющих условный (минимальный) набор продуктов питания Росстата}\label{app:A}

\section{Структура условного (минимального) набора продовольственных товаров Росстата}\label{app:A1}

\begin{longtable}{|p{1.8cm}|p{6cm}|p{4cm}|p{4cm}|} % Ширина столбцов увеличена для занимания всей страницы
	\caption{Перечень продовольственных категорий, составляющих условный (минимальный) набор продуктов питания Росстата} 
	\label{tab:food_categories} \\
	\toprule
	\textbf{№ п/п} & \textbf{Название категории} & \textbf{Единица измерения} & \textbf{Вес (количество) товара в год} \\
	\midrule
	\endfirsthead
	
	\toprule
	\textbf{№ п/п} & \textbf{Название категории} & \textbf{Единица измерения} & \textbf{Вес (количество) товара в год} \\
	\midrule
	\endhead
	
	\hline
	\multicolumn{4}{|r|}{\small \slshape Продолжение следует} \\ 
	\hline
	\endfoot
	
	\endlastfoot
	
	\footnotesize
	1   & Говядина (кроме бескостного мяса) & кг & 15 \\
	2   & Свинина (кроме бескостного мяса) & кг & 4 \\
	3   & Баранина (кроме бескостного мяса) & кг & 1,8 \\
	4   & Куры охлажденные и мороженые & кг & 14 \\
	5   & Рыба мороженая неразделанная & кг & 14 \\
	6   & Сельдь соленая & кг & 0,7 \\
	7   & Масло сливочное & кг & 1,8 \\
	8   & Масло подсолнечное & кг & 7 \\
	9   & Маргарин & кг & 6 \\
	10  & Молоко питьевое цельное пастеризованное 2,5 - 3,2\% жирности & л & 110 \\
	11  & Сметана & кг & 1,8 \\
	
	\end{longtable}
	
	\newpage % Переход на новую страницу
	
	\noindent
	\textbf{Продолжение таблицы \ref{tab:food_categories}:} \\
	
	\begin{longtable}{|p{1.8cm}|p{6cm}|p{4cm}|p{4cm}|} % Уменьшение ширины столбцов
	\toprule
	\textbf{№ п/п} & \textbf{Название категории} & \textbf{Единица измерения} & \textbf{Вес (количество) товара в год} \\
	\midrule
	\endfirsthead
	
	\toprule
	\textbf{№ п/п} & \textbf{Название категории} & \textbf{Единица измерения} & \textbf{Вес (количество) товара в год} \\
	\midrule
	\endhead
	
	\hline
	\multicolumn{4}{|r|}{\small \slshape Продолжение следует} \\ 
	\hline
	\endfoot
	
	\endlastfoot
	
	\footnotesize % Уменьшение размера текста на большее значение
	12  & Творог нежирный & кг & 10 \\
	13  & Сыры сычужные твердые и мягкие & кг & 2,5 \\
	14  & Яйца куриные & 10 шт. & 18 \\
	15  & Сахар-песок & кг & 20 \\
	16  & Мука пшеничная & кг & 20 \\
	17  & Хлеб из ржаной муки и из смеси муки ржаной и пшеничной & кг & 115 \\
	18  & Хлеб и булочные изделия из пшеничной муки 1 и 2 сортов & кг & 75 \\
	19  & Рис шлифованный & кг & 5 \\
	20  & Пшено & кг & 6 \\
	21  & Горох и фасоль & кг & 7,3 \\
	22  & Вермишель & кг & 6 \\
	23  & Картофель & кг & 150 \\
	24  & Капуста белокочанная свежая & кг & 35 \\
	25  & Морковь & кг & 35 \\
	26  & Огурцы свежие & кг & 1,8 \\
	27  & Лук репчатый & кг & 20 \\
	28  & Яблоки & кг & 18,6 \\
	29  & Печенье & кг & 0,7 \\
	30  & Карамель & кг & 0,7 \\
	31  & Чай черный байховый & кг & 0,5 \\
	32  & Соль поваренная пищевая & кг & 3,65 \\
	33  & Перец черный (горошек) & кг & 0,73 \\
\end{longtable}

\noindent
\text{Источник: открытые данные из сети Интернет}


\section{Базовые факты жесткости цен на высокочастотных данных}\label{app:A2}

\begin{longtable}{|p{1cm}|p{8.5cm}|p{3.5cm}|p{3cm}|} % Уменьшение ширины столбцов
	\caption{Средние частоты и дюрации: продовольственные товары} \label{tab:products} \\
	\hline
	\textbf{№ п/п} & \textbf{Наименование категории} & \textbf{Средняя частота (\% в день)} & \textbf{Средняя дюрация (дней)} \\
	\hline
	\hline
	\endfirsthead
	
	\hline
	\textbf{№ п/п} & \textbf{Наименование категории} & \textbf{Средняя частота (\% в день)} & \textbf{Средняя дюрация (дней)} \\
	\hline
	\hline
	\endhead
	
	\hline
	\endfoot
	
	\hline
	\endlastfoot
	
	\footnotesize % Уменьшение размера текста на большее значение
	1   & Баранина (кроме бескостного мяса) & 1,13 & 88,0 \\
	2   & Говядина (кроме бескостного мяса) & 0,76 & 73,6 \\
	3   & Пшено & 1,73 & 57,3 \\
	4   & Рыба мороженая неразделанная & 1,97 & 50,3 \\
	5   & Горох и фасоль & 2,27 & 43,6 \\
	6   & Хлеб и булочные изделия из пшеничной муки 1 и 2 сортов & 2,32 & 42,6 \\
	7   & Свинина (кроме бескостного мяса) & 2,51 & 39,3 \\
	8   & Сахар-песок & 2,63 & 37,5 \\
	9   & Сыры сычужные твердые и мягкие & 2,81 & 35,1 \\
	10  & Рис шлифованный & 2,81 & 35,1 \\
	11  & Карамель & 2,92 & 33,7 \\
	12  & Сельдь соленая & 2,97 & 33,2 \\
	13  & Соль поваренная пищевая & 3,2 & 30,7 \\
	14  & Печенье & 3,29 & 29,9 \\
	15  & Вермишель & 3,59 & 27,4 \\
	16  & Маргарин & 3,59 & 27,4 \\
	17  & Куры охлажденные и мороженые & 3,6 & 27,3 \\
	18  & Крупа гречневая-ядрица & 3,63 & 27,0 \\
	19  & Мука пшеничная & 3,66 & 26,8 \\
	20  & Чай черный байховый & 3,69 & 26,6 \\
	21  & Перец черный (горошек) & 3,71 & 26,5 \\
	22  & Яйца куриные & 3,76 & 26,1 \\
	23  & Молоко питьевое цельное пастеризованное 2,5--3,2\% жирности & 3,86 & 25,4 \\
	24  & Масло сливочное & 3,88 & 25,3 \\
	25  & Творог нежирный & 3,92 & 25,0 \\
	26  & Хлеб из ржаной муки и из смеси муки ржаной и пшеничной & 3,98 & 24,6 \\
	27  & Сметана & 4,03 & 24,3 \\
	28  & Масло подсолнечное & 4,51 & 21,7 \\
	29  & Лук репчатый & 4,56 & 21,4 \\
	30  & Морковь & 4,69 & 20,8 \\
	31  & Капуста белокочанная свежая & 4,98 & 19,6 \\
	32  & Яблоки & 6,64 & 14,6 \\
	33  & Огурцы свежие & 6,75 & 14,3 \\
	34  & Картофель & 7,17 & 13,4 \\
	
\end{longtable}
\noindent
\text{Источник: составлено автором}


\begin{longtable}{|p{1cm}|p{8.5cm}|p{3.5cm}|p{3cm}|} % Задаем ширину столбцов
	\caption{Средние частоты и дюрации: непродовольственные товары}
	\label{tab:non_food_products} \\
	\hline
	\textbf{№ п/п} & \textbf{Наименование категории} & \textbf{Средняя частота (\% в день)} & \textbf{Средняя дюрация (дней)} \\
	\hline
	\hline
	\endfirsthead
	
	\hline
	\textbf{№ п/п} & \textbf{Наименование категории} & \textbf{Средняя частота (\% в день)} & \textbf{Средняя дюрация (дней)} \\
	\hline
	\hline
	\endhead
	
	\hline
	\endfoot
	
	\hline
	\endlastfoot
	
	1 & Джемпер для детей школьного возраста & 0,02 & 4 999,5 \\ \hline
	2 & Кроссовые туфли для детей с верхом из искусственной кожи & 0,06 & 1 666,2 \\ \hline
	3 & Брюки для детей школьного возраста из джинсовой ткани & 0,07 & 1 428,1 \\ \hline
	4 & Метамизол натрия (Анальгин отечественный) & 7,22 & 13,3 \\ \hline
	
\end{longtable}
\noindent
\text{Источник: составлено автором}


\begin{longtable}{|p{1cm}|p{8.5cm}|p{3.5cm}|p{3cm}|} % Ширина столбцов оптимизирована
	\caption{Средние частоты и дюрации: услуги}
	\label{tab:services} \\
	\hline
	\textbf{№ п/п} & \textbf{Наименование категории} & \textbf{Средняя частота (\% в день)} & \textbf{Средняя дюрация (дней)} \\
	\hline
	\hline
	\endfirsthead
	
	\hline
	\textbf{№ п/п} & \textbf{Наименование категории} & \textbf{Средняя частота (\% в день)} & \textbf{Средняя дюрация (дней)} \\
	\hline
	\hline
	\endhead
	
	\hline
	\endfoot
	
	\hline
	\endlastfoot
	
	1 & Постановка набоек & 0 &  \\ \hline
	2 & Стрижка модельная в женском зале & 0,06 & 1 666,2 \\ \hline
	3 & Абонентская плата за неограниченный объем местных телефонных соединений & 0,14 & 713,8 \\ \hline
	4 & Помывка в бане в общем отделении & 0,16 & 624,5 \\ \hline
	5 & Проезд в городском муниципальном автобусе & 0,30 & 332,8 \\ \hline
	6 & Проезд в троллейбусе & 0,32 & 312,0 \\ \hline
	
\end{longtable}
\noindent
\text{Источник: составлено автором}



\section{Сопоставление фактов с предсказаниями моделей ценообразования}\label{app:A3}


\begin{table}[h]
	\centering
	\caption{Сопоставление фактов с предсказаниями моделей ценообразования}
	\label{tab:comparsion}
	\small
	\begin{tabularx}{\textwidth}{|p{4cm}|X|X|X|X|}
		\hline
		Факты, наблюдаемые в данных & \multicolumn{2}{c|}{Модели, зависящие от состояния} & \multicolumn{2}{c|}{Модели, зависящие от времени} \\ \cline{2-5}
		& [Golosov, Lucas, 2007] & [Dotsey, et al., 1999] & [Calvo, 1983] & [Taylor, 1980] \\ \hline
		Относительно большие размеры изменений цен по модулю & Да & Да & Да & Да \\ \hline
		Множество малых (меньше 5\% по модулю) изменений цен & Нет & Да & Да & Да \\ \hline
		Малая доля изменений цен меньше 1\% по модулю - бимодальность в распределении изменений цен & Да & Нет & Нет & Нет \\ \hline
		Вероятность изменения цены увеличивается с длительностью периода «жизни» цены & Да & Нет & Нет & Нет \\ \hline
		Размер изменения цены не увеличивается с длительностью периода «жизни» цены & Да & Нет & Нет & Нет \\ \hline
		Средний размер изменений цен коррелирует сильнее с инфляцией, чем доля меняющихся цен & Нет & Нет & Да & Да \\ \hline
		Дисперсия инфляции в большей степени объясняется дисперсией интенсивной составляющей & Да & Нет & Да & Да \\ \hline
	\end{tabularx}
	
	\vspace{0.5em}
	\footnotesize
	Источник: составлено автором
\end{table}


